% Dueño: B. A17: contexto, trabajos relacionados y estado del arte. A18: termina con la estructura del documento.
\IEEEPARstart{L}{a} esteganografía en señales de audio consiste en ocultar información digital dentro de una señal portadora acústica sin alterar de manera perceptible su reproducción, lo que resulta de gran utilidad en el almacenamiento de metadatos de canciones, el control de derechos de autor y la autenticación de grabaciones~\cite{dutta2009,kaur2014}.

Entre las técnicas clásicas reportadas en el estado del arte destacan la sustitución del bit menos significativo, la codificación de fase, el espectro disperso y el enmascaramiento con eco~\cite{dutta2009,kaur2014}. El enmascaramiento con eco, introducido originalmente por Gruhl, Lu y Bender~\cite{gruhl_echo_hiding}, aprovecha la propiedad psicoacústica según la cual una réplica atenuada con un retardo del orden de un milisegundo no se percibe como un eco separado, sino como una ligera resonancia espectral enmascarada por la señal principal.

El resto de este artículo se organiza de la siguiente manera. En la Sección~II se desarrolla el marco teórico del eco como sistema lineal e invariante en el tiempo (LTI) y del cepstro, y en la Sección~III, la solución propuesta. En la Sección~IV se presentan y discuten los resultados experimentales. Finalmente, en la Sección~V se exponen las conclusiones y las recomendaciones para trabajos futuros.
